\documentclass[conference]{IEEEtran}
\IEEEoverridecommandlockouts
% The preceding line is only needed to identify funding in the first footnote. If that is unneeded, please comment it out.
\usepackage{cite}
\usepackage{amsmath,amssymb,amsfonts}
\usepackage{algorithmic}
\usepackage{graphicx}
\usepackage{textcomp}
\usepackage{xcolor}
\usepackage{booktabs}
\usepackage{multirow}
\usepackage{array}
\usepackage{colortbl}
\def\BibTeX{{\rm B\kern-.05em{\sc i\kern-.025em b}\kern-.08em
    T\kern-.1667em\lower.7ex\hbox{E}\kern-.125emX}}
\begin{document}

\title{Interpretable and Robust Root Cause Analysis for Microservices via Kolmogorov-Arnold Graph Neural Networks}

\author{\IEEEauthorblockN{1\textsuperscript{st} Guan-Hao Huang}
\IEEEauthorblockA{\textit{Department of Computer Science} \\
\textit{National Taichung University of Education}\\
Taichung, Taiwan \\
bcs113116@gm.ntcu.edu.tw}
\and
\IEEEauthorblockN{2\textsuperscript{nd} Yu-Chen Lai}
\IEEEauthorblockA{\textit{Department of Computer Science} \\
\textit{National Taichung University of Education}\\
Taichung, Taiwan \\
bcs113115@gm.ntcu.edu.tw}
\and
\IEEEauthorblockN{3\textsuperscript{rd} Kuan-Chou Lai}
\IEEEauthorblockA{\textit{Department of Computer Science} \\
\textit{National Taichung University of Education}\\
Taichung, Taiwan \\
kclai@mail.ntcu.edu.tw}
}

\maketitle

\begin{abstract}
Despite the exceptional feature extraction capabilities of Graph Neural Networks (GNNs) in microservices root cause analysis, their black-box decision mechanisms and lack of domain inductive bias severely hinder their application in high-stakes critical operations scenarios where system failures may lead to significant losses. To resolve the dilemma between expressiveness and interpretability, this study proposes GNN\_KAN, an innovative architecture based on the Kolmogorov--Arnold representation theorem. GNN\_KAN decomposes multivariate functions into visualizable one-dimensional basis function combinations, endowing the model with intrinsic interpretability while preserving the powerful fitting capabilities of deep learning, providing a new perspective for root cause analysis from implicit feature mapping to explicit function representation. This study further constructs a domain-aware feature engineering framework that systematically eliminates posterior leakage risks in features. Experimental results demonstrate that GNN\_KAN outperforms existing statistical baseline methods (such as BARO) on key metrics including Top-k accuracy and Precision@k across multiple core datasets containing multi-source telemetry data. In complex topology scenarios (Train Ticket), GNN\_KAN achieves Precision@1 of 0.256, representing an 88\% improvement over baseline methods. Particularly in memory fault diagnosis for hyperscale systems, GNN\_KAN achieves Precision@1=0.96, representing over 12-fold improvement compared to baseline methods, confirming its exceptional capability in capturing nonlinear saturation effects. This study further reveals the applicability boundaries between statistical methods and deep learning, and based on extensive empirical research, uncovers the intrinsic mapping relationships between different basis functions and system topology complexity, providing systematic selection strategies for practical applications.
\end{abstract}

\begin{IEEEkeywords}
Root cause analysis, Graph neural networks, Kolmogorov-Arnold Networks, Microservices systems, Explainable AI
\end{IEEEkeywords}

\section{Introduction}

\subsection{Research Background and Motivation}

With the widespread adoption of microservices architectures \cite{microservices_survey}, system topologies have become increasingly complex, and traditional root cause analysis (RCA) faces severe challenges. When system anomalies occur, upstream service anomalies rapidly propagate downstream along dependency relationships, causing multiple downstream services to simultaneously exhibit abnormal symptoms. This ``cascade effect'' often masks the true root cause \cite{train_ticket}. The field of root cause analysis has long faced the \textbf{dilemma between expressiveness and interpretability}: \textbf{statistical and correlation analysis methods} (such as BARO \cite{baro2023}) possess interpretability but lack the ability to fit complex nonlinear propagation; \textbf{deep learning and GNN methods} possess powerful fitting capabilities, but their black-box nature prevents operations personnel from trusting them, and the lack of domain inductive bias limits their generalization ability. Particularly in \textbf{high-stakes scenarios}, i.e., critical operations environments where system failures may lead to significant economic losses, service disruptions, or security risks (such as financial trading systems, healthcare information systems, and e-commerce core services), the interpretability and trustworthiness of model decisions become necessary conditions for deployment, not merely desirable attributes.

Traditional GNNs rely on black-box MLPs (Multi-Layer Perceptrons) for feature transformation, making it difficult for operations personnel to trust their decision paths and verify whether model decisions follow reasonable causal propagation. More critically, activation functions commonly used in MLPs such as ReLU may have limitations when simulating specific physical characteristics of microservices anomaly propagation (such as smooth saturation effects and threshold triggering mechanisms). In contrast, KAN's additive structure and one-dimensional function characteristics provide possibilities for visual verification, enabling operations engineers to directly inspect the function shapes of anomaly propagation. More importantly, KAN, through basis functions such as B-Spline or Chebyshev polynomials, possesses \textbf{smoothness} and \textbf{locality} as \textbf{inductive bias}, making it more aligned with the physical characteristics of real systems when capturing continuously varying anomaly signals compared to MLPs.

\subsection{Research Contributions}

To break the above dilemma, this study proposes \textbf{GNN\_KAN}---an innovative root cause analysis architecture based on the Kolmogorov--Arnold representation theorem. GNN\_KAN deeply integrates KAN into GNN message passing mechanisms, endowing the model with intrinsic interpretability by decomposing high-dimensional functions into visualizable univariate function combinations while preserving the powerful fitting capabilities of deep learning. The main contributions of this study include: (1) proposing a white-box architecture integrating KAN with GNN message passing, enabling operations experts to inspect propagation response curves (such as saturation effects and threshold effects) at the single-step message passing level; (2) constructing an RCA-aware feature engineering framework that eliminates posterior leakage, systematically avoiding future information leakage and annotation coupling risks; (3) based on extensive empirical research, revealing intrinsic mapping relationships between basis functions and system topology complexity (e.g., PQC\_GPU demonstrates 8-fold performance advantage in hyperscale systems), providing systematic selection strategies for practical applications; (4) achieving \textbf{Precision@1=0.96} (baseline only 0.08) for memory faults on the Train-Ticket hyperscale dataset, confirming the model's capability in capturing nonlinear physical characteristics.

The remainder of this paper is organized as follows: Section 2 reviews related work; Section 3 details the methodology; Section 4 describes the experimental setup; Section 5 reports experimental results and analysis; Section 6 discusses research insights and limitations; Section 7 concludes the study and proposes future research directions.

\section{Related Work}

Root cause analysis is a core problem in distributed systems reliability engineering. \textbf{Graph-based statistical methods} (such as MicroRank \cite{microrank2021}, T-Rank \cite{trank2021}, CloudRanger \cite{cloudranger2018}, BARO \cite{baro2023}, PC\_PageRank) utilize random walks, spectral analysis, or PageRank for root cause ranking, possessing intuitive interpretability, but often assume linear propagation, making it difficult to capture nonlinear cascade effects in complex topologies, and difficult to learn edge weights to accurately reflect real propagation strength. \textbf{Deep learning methods} (such as \cite{gnn_rca2023}, Nezha \cite{nezha2023}, EADRO \cite{eadro2023}) improve fitting capabilities, but existing GNN methods mostly use standard MLPs (ReLU/GeLU), lacking interpretability; experimental results show that traditional GNN (MLP) almost fails in complex scenarios (Train-Ticket: P@1=0.024), confirming the limitations of black-box MLPs. \textbf{Causal inference methods} (such as CausalRCA \cite{causalrca2023}) attempt to reconstruct causal graphs, having advantages in interpretability, but face challenges of high computational cost, large sample requirements, and difficult-to-satisfy assumptions in high-frequency, high-dimensional, non-stationary scenarios. \textbf{Knowledge-guided learning} (such as KAN \cite{kan2024}) injects domain knowledge into models. KAN, based on the Kolmogorov--Arnold representation theorem \cite{kolmogorov1957}, decomposes multivariate functions into learnable one-dimensional function combinations, making nonlinear transformations both expressive and interpretable. GNN\_KAN introduces KAN into graph neural network message passing, replacing black-box nonlinearities with KAN, embedding domain knowledge through basis function shapes, balancing expressiveness and interpretability. Recent research such as RCACopilot \cite{rcacopilot2024} and Xpert \cite{xpert2024} explores using LLMs for fault diagnosis, but LLMs face challenges of high inference costs, hallucination problems, and insufficient interpretability when processing high-frequency numerical time-series data, making GNN\_KAN based on numerical computation and visualizable functions still have clear advantages in real-time performance, accuracy, and trustworthiness.

\section{Methodology}

This section first rigorously defines the root cause analysis problem, clarifying inputs, outputs, and main challenges; then elaborates on the overall architecture of GNN\_KAN and how each module collaboratively solves the RCA problem; finally details graph modeling strategies, KAN module design, RCA-aware feature engineering, and learning objectives.

\subsection{Problem Definition and Challenges}

\textbf{Root Cause Analysis (RCA)} aims to identify the most likely root cause that triggers an anomaly from a large number of possible fault points when a distributed system experiences an anomalous event. Given a microservices system graph $G=(V, E)$, where $V$ is the set of service nodes and $E \subseteq V \times V$ is the set of dependency edges. Let $X = \{x_v(t) | v \in V, t \in W\}$ be node observation metrics (such as latency, error rate, throughput, CPU/memory usage) within time window $W$, and $A$ be the anomaly event identifier. The research goal is to learn a ranking function $f: (G, X) \to R$ that computes a root cause score $s_v$ for each node $v$, such that the true root cause node $v^*$ maximizes Top-k accuracy in the ranking list $R$.

\textbf{Core technical challenges} include: (1) \textbf{Cascade masking effect of anomaly propagation}: anomalies from root cause nodes propagate downstream along dependency chains, causing multiple downstream nodes to simultaneously exhibit abnormal symptoms, where downstream nodes' abnormal symptoms may be more obvious, yet the true root cause is often located upstream; (2) \textbf{High-dimensional observations and noise interference}: each node produces multi-dimensional observations (tens to hundreds of dimensions), and observations contain noise, spike interference, periodic fluctuations, and load-related normal variations; (3) \textbf{Dynamic topology and structural uncertainty}: service dependency graphs are not static and may change with version updates, traffic routing adjustments, and dynamic scaling; (4) \textbf{Interpretability and practical trust}: operations engineers need to understand ``why this node is determined as the root cause,'' including supporting evidence, propagation paths, and comparisons with other nodes.

\subsection{GNN\_KAN Solution Architecture}

To address the above challenges, GNN\_KAN proposes an integrated solution achieved through collaborative design of graph neural networks and Kolmogorov--Arnold representation networks. The overall model and processing flow are shown in Fig.~\ref{fig:architecture}:

\begin{figure*}[t!]
    \centering
    \includegraphics[width=0.95\textwidth]{figure1_structure.jpg}
    \caption{\textbf{GNN\_KAN Architecture Overview.} The left side shows the dual-graph structure (similarity graph $G_{\text{sim}}$ and propagation graph $G_{\text{prop}}$) constructed based on RCA-aware features; the middle layer demonstrates the message passing mechanism based on Kolmogorov--Arnold representation, replacing traditional MLPs with learnable one-dimensional functions (wavy line symbols), including three core modules: $\text{KAN}_{\text{edge}}$, $\text{KAN}_{\text{message}}$, and $\text{KAN}_{\text{node}}$; the right side shows the $\text{KAN}_{\text{fusion}}$ fusion layer and readout layer, achieving end-to-end white-box root cause localization. The figure labels two main paths: the similarity path (processing similarity features $h^{sim}$) and the propagation path (processing propagation features $h^{prop}$), ultimately generating root cause scores $s_v$ through the fusion layer.}
    \label{fig:architecture}
\end{figure*}

\textbf{Phase One: Multi-hop Anomaly Propagation Modeling via Graph Message Passing}  
This study represents microservices systems as dynamic weighted directed graphs $G = (V, E, w)$, where $w: E \to \mathbb{R}$ is an edge weight function reflecting anomaly propagation strength. Through multi-layer graph neural network message passing mechanisms, node representations $h_v$ not only integrate their own observations $x_v$ but also aggregate messages from neighbor nodes. This mechanism enables the model to capture multi-hop propagation characteristics: upstream node representations influence downstream nodes after multi-layer aggregation, and the learned edge weights and aggregation weights can reflect the importance of propagation paths. Compared to single-node analysis or pairwise correlation analysis, GNN's structured reasoning can effectively distinguish direct impacts from indirect propagation, mitigating downstream masking effects.

\textbf{Phase Two: Interpretable Nonlinear Modeling via Kolmogorov--Arnold Representation}  
To address the problem of MLPs lacking domain semantics and interpretability in traditional GNNs, this study replaces them with Kolmogorov-Arnold Networks (KAN). KAN, based on the Kolmogorov--Arnold representation theorem, decomposes multivariate functions into combinations of one-dimensional learnable functions. Specifically, KAN constructs function spaces using parameterizable one-dimensional basis functions (such as B-spline, Chebyshev polynomials, Fourier basis, PQC\_GPU), achieving function approximation through learning combination coefficients. This design brings dual advantages: (1) \textbf{Interpretability of function decomposition}---the shape of each one-dimensional function can be directly visualized and verified, conforming to domain intuitions such as ``stronger anomalies lead to more significant propagation''; (2) \textbf{Reduced parameter complexity and overfitting risk}---compared to high-dimensional MLPs, one-dimensional function combinations are superior in parameter efficiency and generalization ability.

\textbf{Phase Three: RCA-aware Feature Engineering Framework}  
Feature design is a critical factor in RCA success or failure. This study proposes ``RCA-aware'' design principles: only use features directly related to root cause localization and stably measurable in practical applications. Specific principles include: (1) \textbf{Forward observability}---only use statistics observable in the forward direction within the event window (such as latency percentiles, error rates, resource utilization), avoiding introducing future information or post-hoc annotation indicators; (2) \textbf{Propagation-oriented}---focus on propagatable signals (such as upstream latency's impact on downstream) and structural relationships (dependency strength, call frequency), excluding features coupled with annotations or not directly related to root causes; (3) \textbf{Consistent standardization}---adopt unified standardization and alignment strategies (such as z-score standardization, temporal alignment), reducing cross-service bias and distribution shift. This framework ensures the model learns general anomaly propagation patterns rather than training-set-specific spurious correlations.

\textbf{Overall Architecture Overview}: Given an anomaly event window, GNN\_KAN first encodes each node's observations into initial representations $h_v^0$; then through $L$ layers of message passing ($L=3$ is a typical configuration), messages on edges in each layer are transformed by KAN\_edge, nodes aggregate neighbor messages, first transform node representations with KAN\_node, then process aggregated messages with KAN\_message and integrate through residual connections; finally, the readout layer maps $h_v^L$ to root cause scores $s_v$ and outputs a ranked list in descending order. During training, labeled root cause windows serve as supervision signals, minimizing ranking loss and regularization terms; during inference, forward propagation is performed on new anomaly events and candidate lists are output.

\subsection{Graph Modeling and Message Passing Mechanism}

This study represents distributed systems within anomaly event windows as dynamic weighted directed graphs $G_t = (V, E, W_t)$, where node set $V = \{v_1, ..., v_n\}$ corresponds to services or components, and edge set $E \subseteq V \times V$ is extracted from service dependency relationships and call tracing information. The edge weight matrix $W_t \in \mathbb{R}^{|V| \times |V|}$ is initially normalized based on inter-service call frequencies (dividing call counts by maximum call counts to ensure weights fall in $[0, 1]$), and adaptively adjusted through learnable parameters during training to reflect actual anomaly propagation strength. This initialization strategy ensures the model starts learning from reasonable structural priors, accelerating convergence and improving robustness.

For node $v_i \in V$, this study converts its time-series data within the observation window into a fixed-dimensional feature vector $\bar{x}_{v_i} \in \mathbb{R}^{d_{\text{feat}}}$ through feature extraction, then obtains initial representation $h_{v_i}^{(0)} \in \mathbb{R}^{d_h}$ through linear transformation and LayerNorm, where $d_h$ is the hidden representation dimension (typical value is 64).

GNN\_KAN captures multi-hop anomaly propagation characteristics through $L$ layers of message passing. At layer $l$ ($l = 1, ..., L$), the following operations are performed sequentially: (1) \textbf{Edge message generation}: For edge $(v_u \to v_v) \in E$, concatenate source node representation $h_u^{(l-1)}$ with edge features $e_{uv}$, then transform through KAN module $\text{KAN}_{\text{edge}}$ to obtain message $m_{u \to v}^{(l)}$; (2) \textbf{Node message aggregation}: Node $v$ receives messages from all in-neighbors $\mathcal{N}_{\text{in}}(v)$, performing weighted aggregation through attention mechanism; (3) \textbf{Node representation update}: First transform node representation $h_v^{(l-1)}$ with $\text{KAN}_{\text{node}}$ to get $\tilde{h}_v^{(l)}$, then process aggregated message $\tilde{m}_v^{(l)}$ with $\text{KAN}_{\text{message}}$ and integrate through residual connection (residual weight $\alpha=0.3$).

The final layer representation $h_v^{(L)}$ is mapped to root cause score $s_v \in (0, 1)$ through the readout layer (linear transformation + sigmoid function). All nodes are ranked by $s_v$ in descending order, outputting candidate list $R$.

\subsection{Kolmogorov-Arnold Networks Module Design}

Traditional neural networks use fixed-shape activation functions (such as ReLU, Sigmoid, Tanh) and multi-layer perceptrons (MLPs). While such black-box nonlinear modules are computationally simple and expressive, they lack interpretability and are prone to overfitting. \textbf{Kolmogorov-Arnold Networks (KAN)} provide a learnable and interpretable function approximation alternative based on the Kolmogorov--Arnold representation theorem.

For input $z = [z_1, \ldots, z_{d_{\text{in}}}]^T \in \mathbb{R}^{d_{\text{in}}}$, the KAN layer based on the Kolmogorov--Arnold representation theorem \cite{kolmogorov1957} is defined \cite{kan2024} as:
\begin{equation}
\phi(z) = \sum_{i=1}^{d_{\text{in}}} \sum_{k=1}^{K} \alpha_{i,k} \cdot b_k(z_i)
\end{equation}
where $K$ is the number of one-dimensional basis functions (typical value is 5-10), $b_k: \mathbb{R} \to \mathbb{R}$ is a one-dimensional basis function, and $\alpha_{i,k} \in \mathbb{R}$ is a learnable coefficient. Compared to traditional MLPs requiring $O(d_{\text{in}} \times d_{\text{out}})$ parameters, KAN only requires $O(d_{\text{in}} \times K)$ parameters (assuming $K << d_{\text{out}}$), reducing parameter complexity by $O(d_{\text{out}}/K)$ times, lowering overfitting risk and improving generalization ability.

This study supports four types of basis functions, each with different characteristics and ``expressiveness--interpretability'' trade-offs: (1) \textbf{Chebyshev polynomial basis} (default): optimal approximation properties and numerical stability, suitable for fan-out topologies; (2) \textbf{B-spline basis}: local support properties, suitable for linear topologies with smooth propagation; (3) \textbf{Fourier basis}: suitable for capturing periodic anomaly patterns; (4) \textbf{Parameterized Quantum Circuit basis (PQC\_GPU)}: a high-expressiveness option for hyperscale, deeply complex topologies, used to handle extreme scenarios where traditional polynomial bases saturate.

\subsection{Training Objectives and Learning Mechanism}

GNN\_KAN's training objective integrates multiple sub-objectives, balancing accuracy, interpretability, and robustness:

The model training objective is to minimize $L_{\text{total}} = L_{\text{rank}} + \lambda_1 L_{\text{edge}} + \lambda_2 L_{\text{kan}}$. Where $L_{\text{rank}} = \sum_{v \neq v^*} \max(0, \text{margin} - (s_{v^*} - s_v))$ is the Margin Ranking Loss, ensuring root cause scores are significantly higher than other nodes; $L_{\text{edge}} = \lambda_1 \sum_{(u,v)\in E} |w_{uv}|$ ($\lambda_1=0.01$) is edge weight sparsification regularization, encouraging learning sparse critical propagation paths; $L_{\text{kan}} = \lambda_2 \sum_{\text{KAN}} \sum_k |\alpha_k - \alpha_{k-1}|$ ($\lambda_2=0.001$) is KAN smoothing regularization, constraining basis coefficient continuity to preserve prior shapes. The Adam optimizer is used with learning rate adopting warm-up + cosine decay schedule, batch size of 32.

\section{Experimental Setup}

\textbf{Datasets}: This study uses five core datasets from the RCAEval platform, covering different scales and fault types: RE1 series (RE1-OB/SS/TT, 375 cases, fan-out/linear/deeply complex topologies), RE2 series (multi-source telemetry data, 270 cases), RE3 series (code-level faults, 90 cases), covering a complete evaluation spectrum from simple to hyperscale, from resource/network to code-level faults. \textbf{Implementation}: All models are implemented based on PyTorch, executed on NVIDIA 2080ti GPU. GNN\_KAN is set to 3 layers (dim=64), using Adam optimizer (lr=0.001, warm-up + cosine decay, batch=32). All methods use the same data split (60\%/20\%/20\%), window length (5 minutes before and after anomaly trigger), and feature extraction pipeline. \textbf{Baselines}: Comparison methods include statistical method BARO, traditional GNN (MLP), and PC\_PageRank. GNN (MLP) uses the same architecture, only replacing KAN with standard MLP (ReLU, hidden layer 64).\footnote{Causal inference methods (such as CausalRCA) are not included in the comparison due to high computational complexity, large sample requirements, and most algorithms assuming DAG while real microservices have cyclic dependencies.}

\section{Experimental Results}

This section reports experimental results of GNN\_KAN and multiple baseline methods, covering comprehensive evaluation across five core datasets (sock-shop-2, RE1, RE2, train-ticket). This study compares GNN\_KAN's performance under different basis functions (Chebyshev, B-spline, Fourier, PQC\_GPU) and contrasts with the following baseline methods: (1) \textbf{BARO}---a statistical method based on Bayesian change point detection; (2) \textbf{GNN (MLP)}---a GNN baseline using traditional MLP, adopting the same RCA-aware feature extraction and graph structure as GNN\_KAN, only replacing the KAN module with standard MLP, to fairly evaluate KAN's relative advantages; (3) \textbf{PC\_PageRank}---a graph analysis method based on PageRank, utilizing service dependency graphs for root cause ranking.

Experimental results systematically reveal the \textbf{applicability boundaries and failure modes} of various methods in RCA tasks from three perspectives: ``black-box MLP,'' ``traditional graph algorithms,'' and ``statistical methods,'' answering the key question of ``why KAN is needed rather than simply MLP or PageRank.'' In simple topology scenarios, BARO still demonstrates advantages of statistical methods, conforming to Occam's razor principle. However, in complex nonlinear scenarios, GNN\_KAN shows irreplaceable strategic value. Most critically, in the Train-Ticket memory fault scenario, GNN\_KAN (PQC) achieves Precision@1=0.96, representing over 12-fold improvement compared to baseline methods (BARO: 0.08, GNN MLP: 0.00), confirming KAN's physical inductive bias's exceptional capability in capturing nonlinear anomaly patterns such as memory saturation and leakage.

\section{Discussion}

\subsection{Qualitative Analysis and Inductive Bias Verification of KAN for RCA}

Through visualization analysis (see Fig.~\ref{fig:kan_viz}), we reveal the essential source of GNN\_KAN's performance advantages from the \textbf{function shape} perspective. KAN-learned basis functions exhibit clear \textbf{``dead zone''} and \textbf{``saturation zone''} structures, corresponding to system resilience under low load and collapse behavior under high load, precisely capturing threshold triggering and saturation behavior that highly match the actual physical characteristics of microservices systems. In contrast, traditional MLPs (using ReLU activation) learn function shapes that present irregular noise distributions on the same propagation paths, failing to exhibit physically meaningful structures. This confirms that GNN\_KAN's advantages stem from \textbf{structured inductive bias} rather than merely increased parameter count.

\begin{figure}[t!]
    \centering
    \includegraphics[width=0.9\columnwidth]{figure2_mlp_vs_kan.jpg}
    \caption{\textbf{Visualization comparison of propagation functions learned by KAN and MLP (qualitative analysis).} The red solid line (KAN, Proposed) exhibits a smooth S-shaped curve conforming to microservices fault propagation physical characteristics, with clear ``Dead Zone'' and ``Physical Saturation'' features, corresponding to system resilience under low load and collapse behavior under high load; the blue dashed line (MLP, Baseline) presents irregular noise fluctuations, lacking interpretable physical structure. This comparison strongly confirms that KAN possesses superior inductive bias, capable of internalizing the essential characteristics of microservices systems.}
    \label{fig:kan_viz}
\end{figure}

\subsection{Research Limitations}

This study has the following limitations: (1) \textbf{Topology dependency}: GNN\_KAN relies on service dependency graphs as structural priors, which may be incomplete or inaccurate in practice (such as missing tracing data, dynamic routing), potentially affecting propagation path capture; (2) \textbf{Annotation bias}: Root cause annotations are generated through manual post-hoc analysis and may contain errors or subjective judgments; (3) \textbf{Computational overhead}: KAN's computational overhead is slightly higher than MLP (inference time increases by approximately 50\%), which may require further optimization in extremely latency-sensitive scenarios; (4) \textbf{Dataset representativeness}: While the microservices benchmarks used in experiments are representative, they may not cover all practical scenarios (such as hyperscale systems, domain-specific applications); (5) \textbf{Absolute performance boundaries}: In hyperscale complex scenarios such as Train-Ticket, even though GNN\_KAN has significant relative improvements compared to baseline methods (P@1=0.256 vs BARO 0.136 vs GNN 0.024), its absolute P@1 is still relatively low, reflecting that precise Top-1 root cause localization itself is an extremely challenging open problem under real multi-hop cascades and incomplete annotations, requiring future combination of human-machine collaboration and Top-k/path-level explanations to improve practical usability.

\section{Conclusion and Future Work}

This study proposes GNN\_KAN, breaking the long-standing ``dilemma between expressiveness and interpretability'' in root cause analysis. Through the Kolmogorov--Arnold representation theorem decomposing multivariate functions into one-dimensional learnable function combinations, GNN\_KAN endows the model with white-box-level interpretability: visualizable ``dead zone/saturation zone'' function shapes precisely correspond to system physical characteristics. RCA-aware feature engineering systematically eliminates posterior leakage, ensuring the model learns general propagation patterns rather than training-set-specific spurious correlations. In comprehensive evaluation across multiple core datasets, GNN\_KAN demonstrates stable advantages compared to statistical methods and traditional GNN (MLP), particularly in hyperscale complex systems, where PQC\_GPU basis functions significantly outperform statistical and graph analysis baselines (Train Ticket: Precision@1=0.256 vs BARO 0.136 vs GNN 0.024 vs PC\_PageRank 0.110). \textbf{Most critically, in the Train-Ticket memory fault scenario, GNN\_KAN (PQC) achieves Precision@1=0.96, representing over 12-fold improvement compared to baseline methods (BARO: 0.08, GNN MLP: 0.00)}, confirming KAN's physical inductive bias's exceptional capability in capturing nonlinear anomaly patterns such as memory saturation and leakage, providing a breakthrough solution for memory anomaly diagnosis in hyperscale cloud-native systems. This study further reveals method applicability boundaries: in simple linear scenarios, statistical methods (BARO) still have advantages, reflecting Occam's razor principle; but in complex nonlinear scenarios, GNN\_KAN demonstrates irreplaceable strategic value. \textbf{GNN\_KAN marks an important step in root cause analysis from black-box toward interpretable, trustworthy AI}, laying the foundation for building human-machine collaborative next-generation intelligent operations platforms. Future work can explore: deep integration of causal inference, standardization of interpretability quantification, AutoML-driven KAN design, heterogeneous graph learning, temporal GNN, few-shot/transfer learning, continual learning, uncertainty quantification, and active learning.

\begin{thebibliography}{99}
\small

\bibitem{kan2024}
Liu, Z., Wang, Y., Vaidya, S., Ruehle, F., Halverson, J., Soljačić, M., Hou, T. Y., \& Tegmark, M. (2024).
\textit{KAN: Kolmogorov-Arnold Networks}.
arXiv preprint arXiv:2404.19756.

\bibitem{baro2023}
Li, Z., et al. (2023).
\textit{BARO: Robust Root Cause Analysis for Microservices via Multivariate Bayesian Online Change Point Detection}.
In \textit{Proceedings of the ACM Web Conference (WWW)}, 2865-2875.

\bibitem{microservices_survey}
Soldani, J., Tamburri, D. A., \& Van Den Heuvel, W. J. (2018).
\textit{The pains and gains of microservices: A Systematic grey literature review}.
\textit{Journal of Systems and Software}, 146, 215-232.

\bibitem{train_ticket}
Zhou, X., et al. (2018).
\textit{Fault Analysis and Debugging of Microservice Systems: Industrial Survey, Benchmark System, and Empirical Study}.
\textit{IEEE Transactions on Software Engineering}, 47(2), 243-260.

\bibitem{kolmogorov1957}
Kolmogorov, A. N. (1957).
\textit{On the representation of continuous functions of several variables by superposition of continuous functions of one variable and addition}.
\textit{Doklady Akademii Nauk SSSR}, 114(5), 953-956.

\bibitem{microrank2021}
Yu, G., et al. (2021).
\textit{MicroRank: End-to-End Latency Issue Localization with Extended Spectrum Analysis in Microservice Environments}.
In \textit{Proceedings of the ACM Web Conference (WWW)}, 3087-3098.

\bibitem{trank2021}
Ye, Z., et al. (2021).
\textit{T-Rank: A Lightweight Spectrum Based Fault Localization Approach for Microservice Systems}.
In \textit{Proceedings of the IEEE/ACM International Symposium on Cluster, Cloud and Grid Computing (CCGRID)}, 61-70.

\bibitem{cloudranger2018}
Wang, P., et al. (2018).
\textit{CloudRanger: Root Cause Identification for Cloud Native Systems}.
In \textit{Proceedings of the IEEE/ACM International Symposium on Cluster, Cloud and Grid Computing (CCGRID)}, 492-502.

\bibitem{gnn_rca2023}
Wang, X., et al. (2023).
\textit{Graph Neural Network-based Root Cause Analysis for Cloud-Native Systems}.
In \textit{Proceedings of the AAAI Conference on Artificial Intelligence}.

\bibitem{nezha2023}
Yu, G., et al. (2023).
\textit{Nezha: Interpretable Fine-Grained Root Causes Analysis for Microservices on Multi-modal Observability Data}.
In \textit{Proceedings of the ACM Joint European Software Engineering Conference and Symposium on the Foundations of Software Engineering (ESEC/FSE)}, 428-440.

\bibitem{eadro2023}
Lee, C., et al. (2023).
\textit{EADRO: An End-to-End Troubleshooting Framework for Microservices on Multi-source Data}.
In \textit{Proceedings of the IEEE/ACM International Conference on Software Engineering (ICSE)}, 756-768.

\bibitem{causalrca2023}
Xin, R., et al. (2023).
\textit{CausalRCA: Causal Inference Based Precise Fine-Grained Root Cause Localization for Microservice Applications}.
\textit{Journal of Systems and Software}, 203, 111724.

\bibitem{rcacopilot2024}
Chen, Z., et al. (2024).
\textit{RCACopilot: Large Language Model Powered Root Cause Analysis Assistance}.
In \textit{Proceedings of the IEEE/ACM International Conference on Software Engineering (ICSE)}, 2123-2135.

\bibitem{xpert2024}
Jiang, Y., et al. (2024).
\textit{Xpert: Empowering Incident Management with Query Recommendations via Large Language Models}.
In \textit{Proceedings of the IEEE/ACM International Conference on Software Engineering (ICSE)}, 1234-1245.

\end{thebibliography}

\end{document}
